\begin{lemma}{Invertierbarkeit des Grenzwerts invertierbarer Operatoren}
             {InvertierbarkeitDesGrenzwertsInvertierbarerOperatoren}
Seien $(E, \norm\cdot_E)$ ein $\K$--Banachraum, $(F, \norm\cdot_F)$ ein
normierter $\K$--Vektorraum, $\folge T n \N \in \Inv(E, F)^\N$ konvergent
und $T \in L(E, F)$ sei der Grenzwert dieser Folge. 
Dann gilt
\begin{enumerate}[(1)]
 \item\label{InvertierbarkeitDesGrenzwertsInvertierbarerOperatoren1}
 $\left(T_n\inv\right)_{n \in \N}$ ist konvergent $\iff$
 $\left(T_n\inv\right)_{n \in \N}$ ist beschränkt.
 \item\label{InvertierbarkeitDesGrenzwertsInvertierbarerOperatoren2}
 Ist eine der äquivalenten Bedingungen aus (1) erfüllt, so ist
 $T$ invertierbar und $T\inv = \Limes n \infty T_n\inv$.
\end{enumerate}
\end{lemma}
\begin{proof} (1) Da konvergente Folgen stets beschränkt sind, ist lediglich
die Umkehrung zu zeigen.\\
Sei also $\left(T_n\inv\right)_{n \in \N}$ beschränkt.\\
Dann gibt es ein $c \in \R_{>0}$ mit
$\forall\, n \in \N: \opnorm{T_n\inv}F E \le c$.\\
Da $(E, \norm\cdot_E)$ vollständig ist, ist $(L(F, E), \opnorm \cdot F E)$ 
ein Banachraum, sodass es ausreicht zu zeigen, dass $\folge {T\inv} n \N$ eine
Cauchyfolge ist.\\
Sei also $\eps \in \R_{>0}$.\\
Da $\folge T n \N$ konvergiert, ist dies eine Cauchyfolge. Somit existiert also
ein $n_0 \in \N$ mit
$\forall\, n, m \in \N_{\ge n_0}: \opnorm{T_n - T_m} E F < \frac \eps {c^2}$.\\
Seien $n, m \in \N_{\ge n_0}$. Dann gilt:
\begin{align*}
\opnorm{T_n\inv - T_m\inv}FE 
 & = \opnorm{\left(T_n\inv T_n \right)(T_n\inv - T_m\inv)}FE\\
   &\le \opnorm{T_n\inv}FE\opnorm{T_nT_n\inv-T_nT_m\inv}F F\\
   &=c\opnorm{T_mT_m\inv -T_nT_m\inv}FF & \quad T_nT_n\inv = T_mT_m\inv\\
   &=c\opnorm{\left(T_m -T_n\right)T_m\inv}FF\\
   & \le c \opnorm{T_m\inv} F E \opnorm{T_m - T_n} E F\\
   & \le c^2 \opnorm{T_m - T_n}EF < \eps
\end{align*}
Damit ist die erste Aussage gezeigt.\absatz
(2) Sei also $\left(T_n\inv\right)_{n \in \N}$ beschränkt. Nach (1) ist diese
Folge dann konvergent gegen ein $T_0 \in L(X, Y)$.\\
Sei $L(F,E) \times L(E, F)$ mit der Produktnorm ausgestattet und
\[\alpha : L(F, E) \times L(E, F) \rightarrow L(E, E), (A,B) \mapsto AB\text,\]
\[\beta  : L(F, E) \times L(E, F) \rightarrow L(F, F), (A,B) \mapsto BA\text.\]
Diese beiden Abbildungen sind offensichtlich bilinear und wegen der
Submultiplikativität der Operatornorm stetig.\\
Dann gilt mit $(1)$: $\left( \alpha(T_n\inv, T_n)\right)_{n \in \N}$ ist
konvergent und
\begin{align*}
\id_E &= \Limes n \infty T_n\inv T_n = \Limes n \infty \alpha (T_n\inv, T_n)
      = \alpha\left( \Limes n \infty T_n\inv, \Limes n \infty T_n\right)
      = \alpha\left(T_0, T\right)\text.
\end{align*}
Damit ist also $T_0$ eine Linksinverse von $T$.\\
Analog:
\begin{align*}
\id_F & = \Limes n  \infty T_n T_n\inv  = \Limes n \infty \beta (T_n\inv, T_n)
      = \beta\left( \Limes n \infty T_n\inv, \Limes n \infty T_n\right)
      = \beta\left(T_0, T\right)\text.
\end{align*}
Daher ist $T_0$ auch eine Rechtsinverse von $T$.\\
Also ist $T$ invertierbar mit $T\inv = T_0$.
\end{proof}